\section{结论与未来工作}
\label{sec:conclusion}

\subsection{结论}
\label{subsec:conclusion_summary}
在本文中，我们提出了一种多平台混合机器人框架，专为复杂生化实验室工作流的端到端自动化量身定制，具体涵盖溶液配置、水滴角测量和用于生物培养的梯度稀释。通过解耦控制架构，我们的系统通过约束感知的经典运动规划保证了宏观移动底盘导航过程中的绝对流体稳定性。对于灵巧的、接触密集型的机械臂操作，我们实施了一种混合学习策略：在Isaac Sim中训练的姿态感知大规模并行深度强化学习（DRL）智能体用于直立器材抓取，以及带有正则化少样本真实世界微调的遥操作引导模仿学习（IL）策略用于精密微量移液和液滴滴加。大量的实证评估表明，我们的框架在所报告操作技能的合并回合中达到82.7\%的成功率，有效地消除了自主实验室中物理试错的高昂成本和风险。

\subsection{局限性与未来工作}
\label{subsec:limitations_future}
尽管取得了令人鼓舞的结果，但仍存在一些局限性。目前，模仿学习的示教获取依赖于手动遥操作，当扩展到数百种不同的化验协议时，这可能成为瓶颈。此外，视觉策略主要针对刚性玻璃药瓶和微孔板进行了优化，需要进一步适应高度可变形的器材或软组织。

所提供的状态同步汇总值满足预设目标，但仍需回合数、原始轨迹、独立参考测量和不确定性区间来核验精度。

在未来的工作中，我们计划将视觉语言模型（VLM）和大语言模型（LLM）集成到机器人系统中。这将使系统能够语义解析高级自然语言指令（例如"准备1:10的梯度稀释系列并测量基底水滴角"），并自动生成分层状态机（HSM）序列和控制基元，从而更接近完全自主的科学发现。
